% RQ2 Architecture: Capture Mitigation Simulation

\subsection{System Overview}

Our simulation framework extends the base DAO governance model with configurable capture mitigation mechanisms. The architecture comprises:

\begin{enumerate}
    \item \textbf{Token Distribution Module}: Generates realistic power-law distributions
    \item \textbf{Mitigation Layer}: Implements vote caps, quadratic, and velocity mechanisms
    \item \textbf{Capture Metrics}: Computes whale influence, Nakamoto coefficient, Gini
    \item \textbf{Governance Engine}: Standard proposal lifecycle with modified vote weights
\end{enumerate}

\subsection{Mitigation Implementation}

\subsubsection{Vote Cap Implementation}

\begin{lstlisting}[caption=Vote cap mechanism (pseudocode)]
def compute_voting_power_capped(agent, cap_percentage):
    raw_power = agent.tokens + agent.delegated_tokens
    cap_value = total_supply * cap_percentage
    return min(raw_power, cap_value)
\end{lstlisting}

\subsubsection{Quadratic Voting Implementation}

\begin{lstlisting}[caption=Quadratic voting mechanism (pseudocode)]
def compute_voting_power_quadratic(agent):
    raw_power = agent.tokens + agent.delegated_tokens
    return sqrt(raw_power)
\end{lstlisting}

\subsubsection{Velocity Penalty Implementation}

\begin{lstlisting}[caption=Velocity penalty mechanism (pseudocode)]
def compute_voting_power_velocity(agent, window, penalty):
    old_tokens = get_tokens_held_before(agent, current_time - window)
    new_tokens = agent.tokens - old_tokens
    return old_tokens + penalty * new_tokens
\end{lstlisting}

\subsection{Whale Modeling}

To test capture resistance, we model whale behavior:

\begin{table}[htbp]
\centering
\caption{Whale agent characteristics}
\label{tab:whale-model}
\begin{tabular}{lcc}
\toprule
Attribute & Whale & Regular \\
\midrule
Token share & 5-15\% each & 0.1-1\% each \\
Participation & 90\% & 30\% \\
Vote alignment & Self-interested & Mixed \\
Sybil capability & Can split addresses & Single address \\
\bottomrule
\end{tabular}
\end{table}

\subsubsection{Sybil Attack Modeling}

Whales may respond to caps by splitting tokens:

\begin{lstlisting}[caption=Sybil attack modeling (pseudocode)]
def whale_sybil_response(whale, cap):
    if whale.tokens > cap and whale.sybil_capable:
        num_addresses = ceil(whale.tokens / cap)
        # Cost: gas for N addresses, coordination overhead
        if benefit(num_addresses) > cost(num_addresses):
            return split_into_addresses(whale, num_addresses)
    return [whale]  # No split
\end{lstlisting}

\subsection{Capture Metric Computation}

Metrics are computed at each simulation step:

\begin{lstlisting}[caption=Capture metrics computation]
def compute_capture_metrics(agents):
    weights = sorted([voting_power(a) for a in agents], reverse=True)
    total = sum(weights)

    # Whale influence (top 10)
    whale_influence = sum(weights[:10]) / total

    # Nakamoto coefficient
    cumsum = 0
    for i, w in enumerate(weights):
        cumsum += w
        if cumsum > total / 2:
            nakamoto = i + 1
            break

    # Gini coefficient
    gini = compute_gini(weights)

    return {
        'whale_influence': whale_influence,
        'nakamoto': nakamoto,
        'gini': gini,
        'capture_risk': whale_influence / nakamoto * (1 + gini)
    }
\end{lstlisting}

\subsection{Governance Throughput Metrics}

We track mitigation impact on governance efficiency:

\begin{description}
    \item[Pass Rate] Fraction of proposals achieving majority
    \item[Quorum Rate] Fraction of proposals achieving quorum
    \item[Time to Decision] Steps from proposal to resolution
    \item[Effective Participation] Voting power actually cast vs. available
\end{description}

\subsection{Simulation Configuration}

Mitigation experiments use extended configuration:

\begin{lstlisting}[caption=Mitigation configuration schema]
mitigation:
  vote_cap:
    enabled: true/false
    cap_percentage: 0.01-0.10
  quadratic:
    enabled: true/false
  velocity:
    enabled: true/false
    window_days: 7-90
    penalty_factor: 0.0-0.5
\end{lstlisting}
